\documentclass{article}
% \usepackage{times}
\usepackage[margin=1.2in]{geometry}
\usepackage[affil-it]{authblk} 
%%Packages
\usepackage{latexsym, amsthm, amscd, amsfonts, mathrsfs, amsmath, amssymb, stmaryrd, tikz-cd, mathrsfs, bbm, url, esint, mathtools, enumerate, caption, subcaption}
\usepackage{dsfont}
\usepackage{blkarray}
\usepackage[utf8]{inputenc} % allow utf-8 input
\usepackage[T1]{fontenc}    % use 8-bit T1 fonts
\usepackage{hyperref}       % hyperlinks
\usepackage{url}            % simple URL typesetting
\usepackage{booktabs}       % professional-quality tables
\usepackage{nicefrac}       % compact symbols for 1/2, etc.
\usepackage{microtype}      % microtypography
\usepackage{xcolor}  
%%algorithm
\usepackage{algorithm}
\usepackage{algpseudocode}
\usepackage{enumitem}
\usepackage{comment}






%%Commands and operators
\newcommand{\tr}{\mathrm{Tr}}
\newcommand{\He}{\mathrm{He}}
\newcommand{\Corr}{\mathrm{corr}}
\newcommand{\cmin}{c_{\min}}
\newcommand{\Cov}{\mathrm{cov}}
\newcommand{\diag}{\mathrm{diag}}
\newcommand{\pr}{\mathbb{P}}
\newcommand{\E}{\mathbb{E}}
\DeclareMathOperator{\CP}{CP}
\DeclareMathOperator{\Var}{Var}
\newcommand{\Prr}{\mathop{\rm Pr}\nolimits}
\DeclareMathOperator{\argmin}{argmin}
\DeclareMathOperator{\argmax}{argmax}
\DeclareMathOperator{\Maj}{Maj}
\DeclareMathOperator{\Ramp}{Ramp}
\DeclareMathOperator{\ReLU}{ReLU}
\DeclareMathOperator{\Unif}{Unif}
\DeclareMathOperator{\Rad}{Rad}
\DeclareMathOperator{\INAL}{INAL}
\DeclareMathOperator{\sign}{sgn}
\DeclareMathOperator{\erf}{erf}
\DeclareMathOperator{\erfc}{erfc}
\DeclareMathOperator{\orb}{orb}
\DeclareMathOperator{\Span}{span}
\DeclareMathOperator{\NS}{NS}
\DeclareMathOperator{\NN}{NN}
\DeclareMathOperator{\SBM}{SBM}
\DeclareMathOperator{\poly}{poly}
\DeclareMathOperator{\TV}{TV}
\DeclareMathOperator{\Inf}{Inf}
\newcommand{\cF}{\mathcal{F}}
\newcommand{\cX}{\mathcal{X}}
\newcommand{\cU}{\mathcal{U}}
\newcommand{\cD}{\mathcal{D}}
\newcommand{\cN}{\mathcal{N}}
\newcommand{\cE}{\mathcal{E}}
\newcommand{\cP}{\mathcal{P}}
\newcommand{\cG}{\mathcal{G}}
\newcommand{\cI}{\mathcal{I}}
\newcommand{\bR}{\mathbb{R}}
\newcommand{\bI}{\mathbb{I}}
\newcommand{\bN}{\mathbb{N}}
\newcommand{\cK}{\mathcal{K}}
\newcommand{\cC}{\mathcal{C}}
\newcommand{\cL}{\mathcal{L}}
\newcommand{\cB}{\mathcal{B}}
\newcommand{\bS}{\mathbb{S}}
\newcommand{\w}{\mathbf{w}}
\newcommand{\x}{\mathbf{x}}
\newcommand{\indep}{\perp\!\!\!\perp}




%%Colors
\definecolor{mydarkblue}{rgb}{0,0.08,0.45}
  \hypersetup{ %
    colorlinks=true,
    linkcolor=mydarkblue,
    citecolor=mydarkblue,
    filecolor=mydarkblue,
    urlcolor=mydarkblue,
    }
\definecolor{DSgray}{cmyk}{0,0,0,0.7}
\definecolor{DSred}{cmyk}{0,0.7,0,0.7}


%%Author notes
\newcommand{\Authornote}[2]{\noindent{\small\textcolor{DSgray}{\sf{
\textcolor{purple}{[#1: #2]}}}}}
\newcommand{\Authornoteo}[2]{\noindent{\small\textcolor{DSgray}{\sf{
\textcolor{orange}{[#1: #2]\marginpar{\textcolor{orange}{\fbox{\Large !}}}}}}}}
\newcommand{\Authornotec}[2]{\noindent{\small\textcolor{DSgray}{\sf{
\textcolor{blue}{[#1: #2]\marginpar{\textcolor{blue}{\fbox{\Large !}}}}}}}}
\newcommand{\enote}{\Authornote{Elisabetta}}
\newcommand{\anote}{\Authornote{Aryo}}


%%Theorem styles
\theoremstyle{plain}
\newtheorem{definition}{Definition}%[section]
\theoremstyle{plain}
\newtheorem{assumption}{Assumption}%[section]
\newtheorem{ass}{Assumption}
\theoremstyle{plain}
\newtheorem{theorem}{Theorem}%[section]
\theoremstyle{plain}
\newtheorem{informaltheorem}{Informal Theorem}%[section]
\theoremstyle{plain}
\newtheorem{exe}{Exercise}%[section]
\theoremstyle{plain}
\newtheorem{example}{Example}%[section]
\theoremstyle{plain}
\newtheorem{fact}{Fact}%[section]
\theoremstyle{plain}
\newtheorem{proposition}{Proposition}%[section]
\theoremstyle{plain}
\newtheorem{claim}{Claim}%[section]
\theoremstyle{plain}
\newtheorem{lemma}{Lemma}%[section]
\theoremstyle{plain}
\newtheorem{corollary}{Corollary}%[thm]
\theoremstyle{plain}
\newtheorem{conjecture}{Conjecture}%[section]
\theoremstyle{remark}
\newtheorem{remark}{Remark}%[section]
\theoremstyle{remark}
\newtheorem{note}{Note}%[section]
% \theoremstyle{discussion}
\newtheorem{discussion}{Discussion}%[section]
\theoremstyle{plain}
\newtheorem{conj}{Conjecture}%[section]
\title{Connectivity}
\author{}
\date{}

\begin{document}

\begin{abstract}
    
\end{abstract}

\maketitle

\section{Introduction}
\begin{itemize}
    \item Motivation of studying how Transformers learn connectivity: Connectivity as a testbed for algorithmic reasoning. References to Ye et al.,2026 and Sanford et al.,2024, Abbe et al., 2024. 
    \item Central question: \textit{How do Transformers learn connectivity? Which architectural features are needed for a given graph structure? }
    \item Previous works study relationship between depth of Transformer and graph diameter for expressivity: Sanford: study depth needed worst cases over graphs and show that depth $O(\log(n))$ is needed, where $n$ is number of nodes. Ye et al. show that depth $ 3^L $ is sufficient and needed, where $L$ is graph diameter. 
    % They consider generalization after training on ER graphs.
    \item We study whether graph diameter is indeed a bottleneck. We restrict to specific graph families: paths, ...  
\end{itemize}

\paragraph{Our Contributions.}

\section{Related Work}

\begin{itemize}
    \item \href{https://arxiv.org/pdf/2405.18512}{Understanding Transformer Reasoning Capabilities
via Graph Algorithms}.
    \item \href{https://arxiv.org/pdf/1307.4884}{Smoothed analysis on connected graphs}: random perturbation of the edge set of an arbitrary (adversarial) graph, makes the diameter to be $O(\log(n))$ with high probability.
\end{itemize}

\section{Setting}
\begin{itemize}
    \item description of architecture.
\end{itemize}


\section{Expressivity with Similarity Readout}

\begin{theorem}
\label{thm:similarity_readout_doubles_radius}
Let $G=(V,E)$ be an undirected graph on $n$ vertices and let
$A\in\{0,1\}^{n\times n}$ be its self-loop-augmented adjacency matrix,
i.e.\ $A_{ij}=1$ if and only if $i=j$ or $\{i,j\}\in E$. Let $d_G$ denote
shortest-path distance in $G$, with $d_G(i,j)=\infty$ if $i,j$ lie in
distinct connected components. Fix $R\geq 1$ and suppose
\[
H(A)=p_R(A)=\sum_{r=0}^{R}\alpha_r A^r,
\qquad \alpha_r\geq 0 \text{ for all } r,\quad \alpha_R>0 .
\]
Let $H_i(A)$ denote the $i$-th row of $H(A)$ and define
\[
s_{ij}(A)=\frac{\langle H_i(A),H_j(A)\rangle}
{\|H_i(A)\|_2\,\|H_j(A)\|_2}.
\]
Then $s_{ij}(A)$ is well defined for all $i,j$, and
\[
s_{ij}(A)>0 \quad\Longleftrightarrow\quad d_G(i,j)\leq 2R .
\]
Consequently, thresholding $s_{ij}(A)$ at $0$ computes all-pairs
connectivity exactly on every graph whose connected components have
diameter at most $2R$.
\end{theorem}


\section{Experimental Results}

\subsection{Path graphs: connectivity can be learned beyond the $2*3^L$ threshold}

\begin{itemize}
    \item odd training. Unbalanced case (exact match perfect), balanced case (exact match not perfect). 
\end{itemize}


\subsection{Difference between paths and cycles}
\begin{itemize}
    \item Motivation: In the attention score, extreme points are important. We compare with cycles, where diameter is the same but no extreme points. 
    \item Paths tend to make far points disconnected, cycles tend to connect disconnected points. 
\end{itemize}

\subsection{Generalization to three paths, finetuning??}


\subsection{Other graph structures: multi-path, barbel}
\begin{itemize}
    \item If somehow helps making the point that diameter is not the right measure. 
    \item If experimental setting is consistent with the previous sections. 
\end{itemize}


\section{Conclusion}


\appendix

\section{Proof of Theorem~\ref{thm:similarity_readout_doubles_radius}}

\begin{proof}
\emph{Step 1: walk characterization of powers of $A$.}
We claim that for every $k\geq 0$ and all $i,j\in V$,
\begin{equation}\label{eq:power_support}
(A^k)_{ij}>0 \quad\Longleftrightarrow\quad d_G(i,j)\leq k .
\end{equation}
The entry $(A^k)_{ij}$ counts walks of length $k$ from $i$ to $j$ in the
augmented graph $G^{\circ}$ obtained from $G$ by adding a self-loop at
every vertex; since all entries of $A$ are nonnegative, $(A^k)_{ij}>0$
if and only if such a walk exists.

($\Leftarrow$) If $d_G(i,j)=\ell\leq k$, take a shortest path
$i=v_0,v_1,\dots,v_\ell=j$ in $G$ and prepend $k-\ell$ traversals of the
self-loop at $i$; this yields a walk of length exactly $k$ in
$G^{\circ}$ from $i$ to $j$.

($\Rightarrow$) Conversely, given a walk of length $k$ from $i$ to $j$
in $G^{\circ}$, delete every self-loop step; the result is a walk in $G$
from $i$ to $j$ of length at most $k$, whence $d_G(i,j)\leq k$.

\emph{Step 2: support of the embeddings.}
Since $\alpha_r\geq 0$ and every $A^r$ is entrywise nonnegative, $H(A)$
is entrywise nonnegative, and
\[
H_{im}(A)=\sum_{r=0}^{R}\alpha_r (A^r)_{im}>0
\quad\Longleftrightarrow\quad
\exists\, r\in\{0,\dots,R\}:\ \alpha_r>0 \text{ and } (A^r)_{im}>0 .
\]
We show this holds if and only if $d_G(i,m)\leq R$. If $H_{im}(A)>0$,
then $(A^r)_{im}>0$ for some $r\leq R$, so by \eqref{eq:power_support}
we get $d_G(i,m)\leq r\leq R$. Conversely, if $d_G(i,m)\leq R$, then
$(A^R)_{im}>0$ by \eqref{eq:power_support}, and since $\alpha_R>0$ the
corresponding term is strictly positive. Hence
\begin{equation}\label{eq:ball_support}
\operatorname{supp}\bigl(H_i(A)\bigr)
= B_R(i) := \{m\in V:\ d_G(i,m)\leq R\}.
\end{equation}
Note that $\alpha_R>0$ alone suffices: no positivity is required of the
lower-order coefficients.

\emph{Step 3: well-posedness.}
For every $i$ we have $d_G(i,i)=0\leq R$, so $i\in B_R(i)$ and
$H_{ii}(A)>0$ by \eqref{eq:ball_support}. Hence $\|H_i(A)\|_2>0$ for all
$i$ and $s_{ij}(A)$ is well defined. Moreover, since the denominator is
strictly positive, $s_{ij}(A)$ and $\langle H_i(A),H_j(A)\rangle$ have
the same sign.

\emph{Step 4: sign of the inner product.}
By nonnegativity of $H(A)$,
\[
\langle H_i(A),H_j(A)\rangle
= \sum_{m\in V} H_{im}(A)\,H_{jm}(A) > 0
\quad\Longleftrightarrow\quad
\exists\, m\in V:\ H_{im}(A)>0 \text{ and } H_{jm}(A)>0,
\]
i.e., by \eqref{eq:ball_support},
\begin{equation}\label{eq:ball_intersect}
\langle H_i(A),H_j(A)\rangle>0
\quad\Longleftrightarrow\quad
B_R(i)\cap B_R(j)\neq\emptyset .
\end{equation}

\emph{Step 5: ball intersection iff distance at most $2R$.}
If $m\in B_R(i)\cap B_R(j)$, then $i$ and $j$ lie in the same connected
component and, by the triangle inequality,
\[
d_G(i,j)\leq d_G(i,m)+d_G(m,j)\leq 2R .
\]
Conversely, suppose $d_G(i,j)=\ell\leq 2R$ and let
$i=v_0,v_1,\dots,v_\ell=j$ be a shortest path. Set
$m=v_{\lceil \ell/2\rceil}$. Along a shortest path,
$d_G(v_0,v_t)=t$ and $d_G(v_t,v_\ell)=\ell-t$ for every $t$, so
\[
d_G(i,m)=\lceil \ell/2\rceil\leq \lceil 2R/2\rceil = R,
\qquad
d_G(m,j)=\ell-\lceil \ell/2\rceil=\lfloor \ell/2\rfloor\leq R,
\]
whence $m\in B_R(i)\cap B_R(j)$. Therefore
\begin{equation}\label{eq:distance_iff}
B_R(i)\cap B_R(j)\neq\emptyset
\quad\Longleftrightarrow\quad
d_G(i,j)\leq 2R .
\end{equation}

\emph{Conclusion.}
Combining Steps~3--5,
\[
s_{ij}(A)>0
\quad\Longleftrightarrow\quad
\langle H_i(A),H_j(A)\rangle>0
\quad\Longleftrightarrow\quad
d_G(i,j)\leq 2R .
\]
If every connected component of $G$ has diameter at most $2R$, then
$d_G(i,j)\leq 2R$ if and only if $d_G(i,j)<\infty$, i.e.\ if and only if
$i$ and $j$ are connected; disconnected pairs have $d_G(i,j)=\infty>2R$
and hence $s_{ij}(A)=0$. Thus $\mathbf{1}\{s_{ij}(A)>0\}$ is exactly the
connectivity matrix of $G$.
\end{proof}

\end{document}
